\section{Conventions}\label{sec:preliminaries}

\paragraph{Weighted games.}
Intermediate instances may have loops, repeated constraint occurrences,
and positive rational weights summing to one. Their value is the maximum
total weight satisfied by a labeling. Occurrences retain separate IDs
even if their constraints coincide. All rational numbers are encoded by
binary numerators and denominators.

A bipartite \emph{projection game} has a question for each vertex on
either side, finite answer sets, and a map from left answers to right
answers on each constraint. A strategy chooses one answer to each
question. Answer sets that depend on the question can be identified
with fixed alphabets of the same size; we use this for the four
satisfying triples of a parity equation. Each prover uses only its own
question. Private or shared randomness independent of the questions
gives a mixture of deterministic strategies and cannot increase value.

The parallel power \(G^{\otimes t}\) samples \(t\) independent
constraint occurrences, sends the ordered question tuples, and accepts
if every coordinate accepts. Answers may depend on the prover's entire
tuple. Repeated question IDs retain separate coordinate positions and
separate local answer copies. This is the ordinary classical product
game used by the repetition theorem.

\paragraph{Binary linear algebra and Fourier analysis.}
Except where a larger field is specified, all vector spaces and maps
are finite-dimensional over \(\F_2\). Write \(\Hom(E,F)\) for the
space of linear maps and \(E^*=\Hom(E,\F_2)\) for the dual. A subspace
\(S\le E^*\) restricts to its image \(S|_K\le K^*\) on a subspace
\(K\le E\); its restriction rank is \(\dim(S|_K)\).

For a real function \(f\) on a binary space \(E\), use normalized
Fourier coefficients
\[
 \widehat f(g)=\E_{x\in E}f(x)(-1)^{g(x)},\qquad
 f(x)=\sum_{g\in E^*}\widehat f(g)(-1)^{g(x)},\qquad
 \sum_g|\widehat f(g)|^2=\E_x|f(x)|^2.
\]
These identities follow because a nonzero linear functional has equally
many zero and one values, so its character has mean zero. The characters
are therefore an orthonormal basis. On a space of matrices, this is
the entrywise pairing with its dual matrix space; we specify tensor
notation when defining the Fourier decoder.

\paragraph{Probability and finite constructions.}
Uniform sampling from a vector space includes zero. Samples are
independent unless a dependence is stated. For finite probability laws,
\[
 \TV(P,Q)=\frac12\sum_x|P(x)-Q(x)|.
\]
It bounds the difference of expectations of any \([0,1]\)-valued
function. We condition only on events of positive probability.

Random experiments describe finite weighted constraint lists.
The reduction enumerates those lists deterministically. Every auxiliary
dimension and tuple length is fixed before the input formula varies;
the final proof checks both this parameter order and the binary
encoding cost. Fourier witnesses and decoders are used only to bound
game value and need not be efficient.
